\documentclass[10pt]{ctexart}
\usepackage[a4paper,margin=1in]{geometry}
\usepackage{amsmath,amssymb,booktabs,array}
\usepackage{enumitem}
\usepackage{listings}
\usepackage{microtype}
\usepackage{hyperref}
\setlist[itemize]{leftmargin=1.5em}
\setlist[enumerate]{leftmargin=1.8em}
\lstset{basicstyle=\ttfamily\small,breaklines=true,columns=fullflexible,keepspaces=true,frame=single,framerule=0.2pt,xleftmargin=0.5em,xrightmargin=0.5em}

\newcommand{\clip}{\operatorname{clip}}
\newcommand{\E}{\mathbb{E}}

\title{016：StepPPO-v3 设计 Review}
\author{}
\date{2026-06-09}

\begin{document}
\maketitle
\vspace{-1.2em}

\section*{结论概览}
StepPPO-v3 的总体方向是合理的。它不再继续调 StepPPO-v2 的混合权重，而是试图把 multi-turn agent 训练中的三个粒度对齐：环境 action 粒度、step-level credit assignment 粒度、以及 PPO clipping 粒度。相对 v2，v3 的核心改动是：删除 raw episode advantage，统一 shaped step reward，使用 pure step GAE advantage，默认 detach value backbone，并把 policy ratio 改成 response/action-level sequence geometric ratio。

该方向能直接回应现有日志里暴露的两个问题：一是 StepPPO-v2 的 reward 口径和 advantage 混合不够干净；二是 GiGPO value-aux detach=false 说明 shared-backbone value regression 本身就可能破坏 policy behavior。因此 v3-core 更像一个必要的诊断版本，而不是额外堆 trick 的复杂版本。

但是，当前 v3 设计在落地前仍有若干需要修正或明确的地方。最重要的是：配置字段需要和现有代码一致，\path{step_ppo_v3} 需要被 dispatch 识别，uid group whitening 的统计口径可能不完全合理，sequence geometric ratio 也不等价于完整 action probability ratio。

\section*{正向评价}
\begin{itemize}
  \item \textbf{reward 口径更干净。} v3 明确要求 policy advantage 和 value target 都来自同一套 shaped step reward，避免 v2 中 episode term 和 step term 分别看到不同 reward 语义。
  \item \textbf{value 干扰被隔离。} 默认 \path{detach_value_backbone=true} 是合理默认值。这样可以先验证 step credit assignment 是否有价值，而不是让 noisy value target 直接牵引 actor backbone。
  \item \textbf{MDP 粒度更一致。} WebShop 环境执行的是完整 textual action；v3 把 response/action 当成一个环境 action，并在 step 轴上做 GAE，比 token-as-step 的解释更自然。
  \item \textbf{diagnostics 更充分。} action ratio、value RMSE/MAE/EV、valid-vs-invalid 分层指标是必要的，否则很难判断失败来自 value、advantage 还是 policy ratio。
  \item \textbf{ablation 顺序合理。} v3-core、no-action-ratio、detach-false、lam-1、batch-norm-adv 这几个对照能较快定位主要因素。
\end{itemize}

\section*{必须先修的工程问题}
\subsection*{1. adv estimator dispatch 不识别 step\_ppo\_v3}
当前实现中 \path{verl/trainer/ppo/step_ppo_algos.py::is_step_ppo_estimator} 只识别字符串尾部等于 \path{step_ppo} 的 estimator：
\begin{lstlisting}
def is_step_ppo_estimator(adv_estimator) -> bool:
    return str(adv_estimator).lower().split(".")[-1] == "step_ppo"
\end{lstlisting}
如果脚本直接设置：
\begin{lstlisting}
algorithm.adv_estimator=step_ppo_v3
\end{lstlisting}
当前 trainer shim 不会调用 StepPPO advantage path。建议改为：
\begin{lstlisting}
def is_step_ppo_estimator(adv_estimator) -> bool:
    name = str(adv_estimator).lower().split(".")[-1]
    return name in {"step_ppo", "step_ppo_v3"}
\end{lstlisting}
同时在 \path{compute_step_ppo_advantage_data} 内根据 estimator 或 config namespace dispatch 到 v2/v3 两套逻辑。

\subsection*{2. policy loss 配置字段写法不匹配}
现有配置结构是：
\begin{lstlisting}
actor_rollout_ref:
  actor:
    policy_loss:
      loss_mode: "vanilla"
\end{lstlisting}
现有 actor 代码读取的是：
\begin{lstlisting}
loss_mode = self.config.policy_loss.get("loss_mode", "vanilla")
\end{lstlisting}
因此 v3 脚本中不应写：
\begin{lstlisting}
actor_rollout_ref.actor.policy_loss=action_level
\end{lstlisting}
这会把 dict 覆盖成字符串，并导致 \path{.get} 失败。应改为：
\begin{lstlisting}
actor_rollout_ref.actor.policy_loss.loss_mode=action_level
\end{lstlisting}
并在 \path{StepWisePPOActor.update_policy} 中显式支持 \path{action_level}。

\subsection*{3. 不能重复扣 invalid action penalty}
当前主训练循环已经在 \path{apply_invalid_action_penalty} 中修改 \path{token_level_scores}，并在已有 \path{step_rewards} 时同步扣除 step reward。v3 文档中的 \path{build_shaped_step_rewards} 如果再次执行：
\begin{lstlisting}
rewards -= invalid_action_penalty_coef * (1 - is_action_valid)
\end{lstlisting}
就会双重惩罚 invalid action。

建议把 reward source 明确分成两种：
\begin{enumerate}
  \item 如果 batch 已经存在 \path{step_rewards} 且 invalid penalty 已由 trainer 扣过，则 v3 直接读取该字段。
  \item 如果需要在 v3 helper 中构造 shaped reward，则 trainer 侧不要提前扣，或者 helper 读取 raw reward 字段并只扣一次。
\end{enumerate}
最稳妥的实现是新增字段 \path{raw_step_rewards} 和 \path{value_step_rewards}，避免同名字段被不同阶段原地修改后语义不清。

\subsection*{4. v3 需要保留 scalar step advantage 字段}
现有 StepPPO 代码会把 final step advantage broadcast 成 token-shaped \path{advantages}，同时保留 scalar \path{step_advantages}。v3 如果要实现 action-level loss，应该直接读取 scalar \path{step_advantages}；token-shaped \path{advantages} 只能作为兼容 fallback。

建议 actor update path 增加：
\begin{lstlisting}
if loss_mode == "action_level":
    step_advantages = micro_data["step_advantages"]
    pg_loss, metrics = compute_policy_loss_action_level(...)
else:
    advantages = micro_data["advantages"]
\end{lstlisting}
这样可以避免 action-level loss 实际仍然通过 token broadcast 和 token aggregation 近似实现，导致 metric 语义不清。

\section*{主要算法风险}
\subsection*{1. uid whitening 可能混合不可比状态}
v3 计划在同一 \path{uid} 内对所有 step samples 做 whitening。这个选择继承了 GiGPO/GRPO 的 group-relative 稳定性，但在 step-level setting 中存在新问题：同一个 task 的 step 0、step 1、step 2 状态分布并不相同，早期探索动作和后期购买动作也不一定可比。

如果把同一个 uid 下所有 step 的 raw GAE advantage 一起 whitening，可能出现两个副作用：
\begin{itemize}
  \item 早期必要但 reward 延迟的动作被后期高 reward step 压成负 advantage。
  \item 不同 trajectory 长度的 step 数不同，长 trajectory 对 uid statistics 贡献更多。
\end{itemize}
建议首轮至少记录 step-index 分层指标：
\begin{lstlisting}
critic/step_adv_norm_mean_by_step_idx
critic/no_variance_group_ratio_by_step_idx
actor/action_clipfrac_by_step_idx
\end{lstlisting}
并把 ablation 从单纯 \path{batch} norm 扩展到：
\begin{lstlisting}
advantage_norm in {uid, uid_step_idx, batch}
\end{lstlisting}
其中 \path{uid_step_idx} 表示同一 uid 且同一 step index 内做 normalization。

\subsection*{2. gamma=0.95, lambda=0.90 可能传播太短}
v3 默认 \(\gamma=0.95\)、\(\lambda=0.90\)，有效 GAE 传播系数是：
\[
\gamma\lambda = 0.855.
\]
如果 WebShop reward 主要在 terminal 或 late-stage 出现，这个设置可能显著削弱远期 credit。降低方差是好事，但不能让 early-step credit 几乎消失。

建议 v3-core 的第一轮不要只跑 \(\lambda=0.90\)，而是至少保留：
\[
\lambda \in \{0.90, 1.00\}
\]
如果 \(\lambda=1.0\) 明显改善 valid action 或 validation，再考虑细调 \(0.95\)。必要时也应测试 \(\gamma=1.0\)，因为当前 WebShop episode 长度较短，\(\gamma=1.0\) 的方差成本可能可控。

\subsection*{3. sequence geometric ratio 不是完整 action probability ratio}
v3 使用：
\[
\rho_i=\exp\left(\frac{1}{L_i}\sum_l m_{i,l}\Delta\log\pi_{i,l}\right).
\]
它比 raw sequence ratio 更稳定，但它约束的是平均 per-token log-ratio，不是完整 response probability ratio：
\[
\rho_i^{\mathrm{sum}}=\exp\left(\sum_l m_{i,l}\Delta\log\pi_{i,l}\right).
\]
因此长 response 的总概率变化仍可能很大，只是平均变化不大。这个设计可以作为工程折中，但文档里应明确它是 trust-region proxy，而不是严格 action probability ratio。

建议额外记录：
\begin{lstlisting}
actor/action_log_ratio_sum_mean
actor/action_log_ratio_mean_mean
actor/action_ratio_geomean_mean
actor/action_ratio_sum_clipped_fraction
\end{lstlisting}
其中 sum ratio 不一定用于 loss，但用于发现长 response drift。

\subsection*{4. action-level clipping 可能过度保守或降低梯度强度}
相比 token-level PPO，action-level geometric ratio 会让整个 response 共用一个 clip decision。优点是粒度一致；风险是某几个 token 的变化可能导致整段 action 被 clip，或者平均后掩盖局部 format token 的变化。

因此 \path{v3_no_action_ratio} 是必要 ablation。建议成功判据不要只看 validation score，还要比较：
\begin{itemize}
  \item valid action ratio；
  \item response length clip ratio；
  \item action clipfrac 与 token pg clipfrac；
  \item old/new rollout probability diff；
  \item action format 错误样例。
\end{itemize}

\subsection*{5. terminal/truncation 语义需要明确}
现有 \path{compute_step_gae_advantage_return} 对每条 trajectory 的最后一个 collected step 直接设 \path{nonterminal=0}。如果最后一步是环境自然 done，这是合理的；如果是 \path{max_steps} 截断，则严格来说应该 bootstrapping 到 next state value，或者至少记录为 truncated case。

v3 应在 rollout 中保存：
\begin{lstlisting}
done_i
truncated_i
\end{lstlisting}
GAE 中使用：
\[
\delta_i = r_i + \gamma (1-done_i) V(s_{i+1}) - V(s_i),
\]
并单独分析 truncation 比例。如果 WebShop 很多轨迹被 max step 截断，当前最后一步 return 会系统性偏低。

\section*{建议的实现顺序}
\begin{enumerate}
  \item 扩展 \path{is_step_ppo_estimator}，让 \path{step_ppo_v3} 被 trainer shim 正确接管。
  \item 新增 \path{algorithm.step_ppo_v3} namespace，但不要复用 v2 的 \path{algorithm.step_ppo} 字段。
  \item 明确 shaped reward 只扣一次 invalid penalty，保留 \path{raw_step_rewards} 和 \path{value_step_rewards}。
  \item 新增 \path{compute_step_ppo_v3_advantage_return}，输出 scalar \path{step_advantages}、\path{step_returns} 和 diagnostics。
  \item 新增 \path{compute_policy_loss_action_level}，不要只 alias 到 GSPO；至少要返回 action-level metrics。
  \item 修改 \path{StepWisePPOActor.update_policy}，支持 \path{policy_loss.loss_mode=action_level} 和 scalar advantage。
  \item 加 2GPU smoke，先检查 metric 是否非 NaN，再跑 4GPU full。
\end{enumerate}

\section*{建议的 v3-core 配置修正版}
\begin{lstlisting}
algorithm.adv_estimator=step_ppo_v3
algorithm.gamma=0.95
algorithm.step_ppo_v3.reward_source=shaped_step
algorithm.step_ppo_v3.use_episode_advantage=false
algorithm.step_ppo_v3.step_gamma=0.95
algorithm.step_ppo_v3.step_lam=0.90
algorithm.step_ppo_v3.advantage_norm=uid
algorithm.step_ppo_v3.zero_no_variance_group=true
algorithm.step_ppo_v3.min_group_std=1.0e-6
algorithm.step_ppo_v3.policy_ratio=sequence_geometric
algorithm.step_ppo_v3.final_advantage_mode=pure_step
actor_rollout_ref.actor.policy_loss.loss_mode=action_level
actor_rollout_ref.actor.loss_agg_mode=seq-mean-token-mean
actor_rollout_ref.actor.value_head.detach_value_backbone=true
actor_rollout_ref.actor.value_head.value_loss_coef=0.03
actor_rollout_ref.actor.value_head.cliprange_value=0.5
\end{lstlisting}

\section*{Smoke Test Checklist}
2GPU smoke 不建议直接看 final score，应先确认以下指标：
\begin{enumerate}
  \item \path{actor/action_ratio_mean}、\path{actor/action_clipfrac}、\path{actor/action_approx_kl} 正常出现。
  \item \path{actor/value_rmse}、\path{actor/value_mae}、\path{actor/value_explained_variance} 非 NaN。
  \item detach=true 下，\path{actor/grad_norm} 不因 value loss 明显放大。
  \item \path{episode/valid_action_ratio} 不长期低于 0.5。
  \item \path{response_length/clip_ratio} 不快速升到 StepPPO-v2 后期水平。
  \item \path{critic/no_variance_group_ratio} 不长期接近 1。
  \item invalid vs valid action 的 advantage/value loss 分层符合预期。
\end{enumerate}

\section*{最终判断}
StepPPO-v3 是一个值得实现的清晰诊断版本，但不应把它直接视为必然更强的算法。它最可能验证的命题是：当 reward 口径一致、value loss 不干扰 actor backbone、policy ratio 与 environment action 粒度一致时，step-level GAE 是否真的能提升 WebShop multi-turn credit assignment。

若 v3-core 仍明显弱于 GiGPO，优先排查三件事：\path{uid} whitening 是否消掉有效信号、\(\gamma\lambda\) 是否传播太短、action-level geometric ratio 是否过度保守。不要在这些问题未定位前引入更复杂 critic 或更多 auxiliary loss。

\end{document}
